\subsection{Learning Without Rational Expectations}
\label{section:fc_paradigms}

This chapter studies world models in reinforcement learning, learned forecasters of state transitions and rewards that an agent uses to plan its actions. Economics has studied the same object since at least the 1980s under the heading of learning without rational expectations, where an agent is not handed the true law of motion of its environment but estimates one from data and acts on the estimate. Write the agent's model at date $t$ as a parameter $\theta_t \in \Theta$ indexing a forecaster $\widehat{P}_{\theta_t}(s' \mid s, a)$ of transitions and $\widehat{r}_{\theta_t}(s, a)$ of rewards. A planning operator $G$ returns a policy $\pi_t = G(\theta_t)$ that is optimal, or at least improved, under the current model, the agent draws a real transition $(s_t, a_t, r_t, s_{t+1})$ with $a_t \sim \pi_t(\cdot \mid s_t)$, and a learning operator $F$ revises the model,
\[
  \theta_{t+1} = \theta_t + \alpha_t \, F(\theta_t, s_t, a_t, r_t, s_{t+1}),
\]
with gain sequence $\alpha_t$, in the form studied in stochastic approximation. A rest point is a $\theta^\star$ at which the model is self-consistent with the data its own induced policy $G(\theta^\star)$ generates, and two questions decide whether the learner succeeds. The first is whether $\theta_t$ converges to a rest point. The second is whether the policy $G(\theta^\star)$ there is optimal on the true environment, which it need not be when the model class cannot contain the truth. The two economic instances below answer one question each, and the reinforcement-learning line that occupies the rest of the chapter is the same loop pursued at scale.

\subsubsection{Two economic instances of model-based learning}
\label{section:fc_paradigm_family}

The first instance is recursive least squares adaptive learning \citep{MarcetSargent1989, EvansHonkapohja2001}. The agent holds a perceived law of motion, a parametric forecast of the variables it cares about, and updates its coefficients $\theta$ from realized data. With forecast target $p_t$, regressor $z_{t-1}$, and perceived disturbance $\varepsilon_t$, the perceived law of motion is $p_t = \theta' z_{t-1} + \varepsilon_t$ and the update is recursive least squares,
\[
  \theta_t = \theta_{t-1} + t^{-1} R_t^{-1} z_{t-1}\big(p_t - \theta_{t-1}' z_{t-1}\big), \qquad
  R_t = R_{t-1} + t^{-1}\big(z_{t-1} z_{t-1}' - R_{t-1}\big),
\]
with $R_t$ the running second-moment matrix of the regressor and $t^{-1}$ the decreasing gain. Decisions taken under the current belief feed back into the data, so a believed $\theta$ induces an actual law of motion whose coefficients form a map $T(\theta)$. A rational-expectations equilibrium is a fixed point $\bar\theta = T(\bar\theta)$, and convergence to it is governed by expectational stability, the local stability of $d\theta/d\tau = T(\theta) - \theta$ at $\bar\theta$, equivalently that the eigenvalues of $DT(\bar\theta) - I$ have strictly negative real parts. In a market with a single feedback coefficient $\lambda$ from average expectations to the realized outcome this reduces to $\lambda < 1$. This is the benign case, since the forecaster can contain the truth and the only thing that can fail is dynamic, with the lone eigenvalue condition deciding whether the belief is drawn to the equilibrium or driven away.

The second instance is misspecified Bayesian learning, formalized as Berk-Nash equilibrium \citep{EspondaPouzo2016}, where the model state is a belief $\mu_t \in \Delta(\Theta)$ rather than the point estimate $\theta_t$ of the loop above. The agent carries a subjective class of models $\{Q_\theta : \theta \in \Theta\}$ of the map from action $a$ to consequence $Y$, with payoff $u(a, Y)$ and reward $r = u(a, Y)$, and forms the predictive $\bar{Q}_\mu(\cdot \mid s, a) = \int_\Theta Q_\theta(\cdot \mid s, a)\, \mu(d\theta)$. Under repeated interaction the belief concentrates on the parameters minimizing a weighted Kullback-Leibler divergence to the true distribution $Q_\pi$ generated under the agent's own policy $\pi$,
\[
  K(\pi, \theta) = \sum_{s, a} \pi(a \mid s)\, \rho(s)\; \mathbb{E}_{Q_\pi(\cdot \mid s, a)}\!\left[\ln \frac{Q_\pi(Y \mid s, a)}{Q_\theta(Y \mid s, a)}\right], \qquad
  \Theta(\pi) = \arg\min_{\theta \in \Theta} K(\pi, \theta),
\]
where $\rho(s)$ is the distribution of the signal $s$. A Berk-Nash equilibrium is a pair $(\pi, \mu)$ with the policy optimal under the predictive belief, so that $\pi(a \mid s) > 0$ implies $a \in \arg\max_{a'} \mathbb{E}_{\bar{Q}_\mu(\cdot \mid s, a')}[u(a', Y)]$, and the belief supported on the best-fit set, $\operatorname{supp}(\mu) \subseteq \Theta(\pi)$. The weighting by the agent's own policy is decisive, since the agent only observes data on the paths its behavior visits. This is the adversarial case, and the ordinary one once the model class of a neural network or an ensemble cannot contain the true dynamics, since the agent converges calmly to a wrong model and a suboptimal policy that the data it generates never contradicts and more data does not dislodge, which is the economic statement of the world-model failure the synthesis returns to in \S\ref{section:fc_synthesis_fails}. The well-specified counterpart in reinforcement learning is posterior sampling, which draws a model from the posterior and plans against it.

Fictitious play is the same loop in the multi-agent case, the estimated model being the empirical distribution of opponents' actions and the planner a best response \citep{Brown1951, FudenbergKreps1993, HofbauerSandholm2002, FudenbergLevine1998, Young2004}.

\label{section:fc_convergence_question}%
Not every economic learning scheme holds a forecaster of the environment. Genetic-algorithm learning \citep{Arifovic1994, Arifovic1995} evolves a population of decision rules directly by realized fitness, the economic counterpart of gradient-free policy search, and evolutionary heuristic switching \citep{BrockHommes1997} selects among a fixed menu of forecasting rules by a logit on past performance, learning selection weights rather than a model. Both omit the learning operator $F$ and mark the model-free boundary the two instances above are defined against. The shared vocabulary of stochastic approximation and fixed-point analysis traces back to \S\ref{section:rl_algorithms}, and all of these modes are run head to head on a fixed environment, one mode at a time, in the dual simulation of \S\ref{section:fc_dual_sim}.

\subsection{Two 1990 Origins: Dyna and Differentiable World Models}
\label{section:fc_origins1990}

Dyna and the controller-model architecture are the reinforcement-learning instances of the model-based loop set out above, with the model-free schemes as their foil. Two papers published in 1990 introduced what the modern literature now calls a world model. In each, an agent learns a forecaster of its environment from interaction data and uses simulated experience drawn from that forecaster to improve its policy, rather than waiting for real transitions to accumulate. \citet{sutton1990} arrived at the architecture through dynamic programming and reinforcement learning, framing it as the unification of direct and indirect updates within a single control loop. \citet{Schmidhuber1990} arrived at it through recurrent neural networks, framing it as a way to make the environment itself differentiable for credit assignment. The two papers do not cite one another and use different vocabularies, but the architectural commitment is the same. We treat them in parallel rather than in lineage, since each anticipates a distinct branch of the modern literature.

\subsubsection{Sutton's integrated architecture}
\label{section:fc_origins_sutton}

\citet{sutton1990} proposes Dyna as an architecture that interleaves three processes at every real time step. The agent performs a direct Q-learning update from the observed transition $(s, a, r, s')$, then updates a learned model $\widehat{P}(s' \mid s, a)$ and $\widehat{r}(s, a)$ from the same transition, then performs $K$ planning updates by sampling previously visited state-action pairs $(\tilde{s}, \tilde{a})$, querying the model for simulated outcomes $(\tilde{r}, \tilde{s}')$, and applying the same Q-update to those synthetic transitions. The ratio $K$ of planning steps to real steps is the lever that controls how heavily the agent exploits its learned model between actions, and is the first explicit statement of the sample-efficiency claim that modern model-based work continues to revisit; the MBPO short-rollout analysis in \citet{janner2019model} is in the same family. We defer the algorithm block to \S\ref{section:fc_dyna_q}, where we develop Dyna-Q in the notation of this survey.

The architectural openness of Dyna is as important as its sample-efficiency claim. Sutton states explicitly that the learned model can be probabilistic and oftentimes incorrect, with the planner doing useful work as long as the cost of synthetic experience is small relative to the cost of real experience. No structure is imposed on the class of admissible models, on the loss used to fit them, or on the planner that consumes their output. This is the architectural slot that subsequent work fills with tabular counts, Gaussian processes, deep ensembles, recurrent state-space models, and implicit latent dynamics, each plugged into the same outer loop. This open slot is what makes Dyna the entry point for the line that runs through \S\ref{section:fc_mbpo} and \S\ref{section:fc_tdmpc2}.

\subsubsection{Schmidhuber's controller-model architecture}
\label{section:fc_origins_schmidhuber}

\citet{Schmidhuber1990} proposes a paired architecture consisting of a recurrent model network $M$ and a controller network $C$. The model $M$ receives the controller's outputs together with current observations and is trained, self-supervised, to predict the next inputs to $C$, including sensory observations and reinforcement signals. The controller $C$ produces actions and is trained jointly with $M$. Once $M$ is a reasonable forecaster, the environment becomes differentiable from the controller's perspective, since the Jacobian of $M$ stands in for the unknown Jacobian of the true dynamics. Planning is then implemented as gradient descent through the model network, by unrolling $C$ and $M$ together for a fixed horizon and descending on a predicted pain or cost signal, without executing any real actions during the planning phase. This is the first proposal of analytic gradient backpropagation through a learned dynamics model for policy improvement.

In the notation of this survey, let $M_\theta$ parameterize the world model with $\widehat{P}_\theta(s_{t+1} \mid s_t, a_t)$ and let $C_\phi$ parameterize the controller with $\pi_\phi(a_t \mid s_t)$. The model is fit on observed transitions by the negative log-likelihood
\[
  \mathcal{L}_M(\theta) = -\sum_t \log \widehat{P}_\theta(s_{t+1} \mid s_t, a_t),
\]
and the controller is fit by maximizing the expected discounted return under the \emph{learned} model,
\[
  J(\phi) = \mathbb{E}_{a_t \sim \pi_\phi(\cdot \mid s_t),\; s_{t+1} \sim \widehat{P}_\theta(\cdot \mid s_t, a_t)} \Big[ \sum_t \gamma^t r_t \Big],
\]
with the planning operator $\phi \leftarrow \phi + \eta \, \nabla_\phi J(\phi)$ implemented by backpropagating through the $H$-step unrolled $M_\theta$ via reparameterized or score-function gradients. A curiosity term adds an intrinsic reward proportional to the model's own prediction error, training $M_\theta$ to reduce it and drawing $C_\phi$ toward states where the model is uncertain. The idea reappears as the intrinsic curiosity of \citet{Pathak2017}, and the full deep realization of the controller-model architecture is \citet{HaSchmidhuber2018}, treated in \S\ref{section:fc_ha_schmidhuber}.

Dyna and the controller-model architecture differ in two places and agree in a third. The model class differs, a tabular Dirac in Dyna against a differentiable network in Schmidhuber, and the planner differs, $K$ tabular $Q$-updates per real step against a gradient step through the unrolled $M_\theta$, but the outer loop is identical, since both learn $M$ from real data, use it to improve the policy, and act in the world. Each anticipated, without the tools to make it work at the scale of pixels or continuous control, a component that took three decades to mature, intrinsic curiosity on one side and planning amplification through short model rollouts on the other.
